\chapter{Introduction}

\section{Background and Motivation}

The Earth's lithosphere, comprising the crust and the uppermost rigid mantle, is a complex and dynamic system that holds the key to understanding the tectonic evolution of our planet. Deciphering the thermal, compositional, and structural state of the lithosphere requires the integration of various geophysical techniques. Among these, the Magnetotelluric (MT) method has emerged as a profoundly powerful tool. Magnetotellurics is a passive, surface-based electromagnetic geophysical technique that measures naturally occurring, time-varying electric and magnetic fields at the Earth's surface to determine the electrical resistivity distribution of the subsurface. Because electrical resistivity is highly sensitive to variations in temperature, partial melt, volatile content (such as water or carbon dioxide), and the presence of interconnected conductive minerals (like graphite or sulfides), MT provides unique constraints on the physicochemical state of the deep Earth that are often complementary to seismic, gravity, and magnetic data.

Over the past several decades, the application of MT has seen exponential growth, transitioning from simple 1D approximations to complex 3D modeling of massive regional arrays. This evolution has been largely driven by advancements in computational power, instrumentation, and numerical algorithms. In the context of cratonic environments and surrounding mobile belts, MT has been instrumental in delineating the thickness of the lithosphere, often referred to as the lithosphere-asthenosphere boundary (LAB), and in identifying fossilized subduction zones, fluid pathways, and deep-seated structural architectures. Southern Africa, the geographical focus of this research, represents an unparalleled natural laboratory for such studies. It hosts some of the Earth's oldest preserved continental crust, including the Archean Kaapvaal and Zimbabwe cratons, bordered by a complex mosaic of Proterozoic mobile belts such as the Limpopo, Namaqua-Natal, and Damara orogens. Understanding the deep structure of these tectonic units is essential for unravelling the processes of early continent formation, stabilization, and subsequent modification.

The Southern African Magnetotelluric Experiment (SAMTEX) was a monumental, international, collaborative effort initiated in the early 2000s, designed to image the electrical structure of the lithosphere across this diverse tectonic collage. The SAMTEX project accumulated an unprecedented volume of broadband and long-period MT data, encompassing thousands of kilometers of transects. While previous studies have utilized subsets of the SAMTEX data to build deterministic 2D and 3D resistivity models, these conventional approaches possess inherent limitations. Deterministic inversion methods typically rely on regularized, gradient-based optimization algorithms (e.g., Occam's inversion) that seek a single "best-fit" model which minimizes a predefined objective function comprising data misfit and a model roughness penalty. Consequently, the resulting models are heavily influenced by the choice of regularization parameters, the starting model, and the inherent non-uniqueness of the electromagnetic inverse problem. Crucially, deterministic methods struggle to provide comprehensive, quantitative assessments of model uncertainty, leaving interpretations potentially vulnerable to artefacts or over-fitting.

The primary motivation for this thesis is to address these limitations by re-evaluating a carefully curated subset of the SAMTEX dataset using an advanced, fully probabilistic approach: Trans-dimensional Hierarchical Bayesian Inversion. Unlike deterministic algorithms, Bayesian inference does not seek a single optimal model; instead, it characterizes the entire ensemble of acceptable models that fit the data within its uncertainties. This is formally expressed as the posterior probability density function (PDF). By employing Reversible Jump Markov Chain Monte Carlo (RJMCMC) algorithms, the inversion becomes "trans-dimensional," meaning the parameterization of the model itself—specifically, the number of subsurface layers in a 1D approximation—is treated as an unknown variable to be solved for. This allows the data to naturally dictate the required complexity of the geoelectric structure, adhering to the principle of parsimony. Furthermore, the incorporation of a hierarchical framework allows for the autonomous estimation of data noise, mitigating the subjective biases often introduced during manual data weighting. By applying this rigorous statistical methodology to diverse tectonic environments in Southern Africa, this research aims to provide a more objective, uncertainty-quantified understanding of the lithospheric architecture, paving the way for more robust tectonic interpretations.

\section\label{sec:data_sources}
\section{Key Data Sources and Tools}

The success of any geophysical investigation is fundamentally tethered to the quality of the observed data and the sophistication of the analytical tools employed. This research leverages a vast, publicly accessible repository of electromagnetic data and an array of modern, open-source computational libraries alongside custom-developed numerical algorithms.

\subsection{Data Sources}
The primary dataset utilized in this dissertation originates from the Southern African Magnetotelluric Experiment (SAMTEX). The SAMTEX array represents one of the largest and most comprehensive regional MT surveys ever conducted, involving the deployment of hundreds of broadband and long-period MT stations across South Africa, Botswana, and Namibia. The overarching goal of SAMTEX was to probe the deep lithospheric roots of the Archean cratons and the surrounding Proterozoic terrains, assessing variations in lithospheric thickness and identifying major structural boundaries.

For the purpose of this study, the processed MT impedance tensor and apparent resistivity data were acquired from the Incorporated Research Institutions for Seismology (IRIS) Data Management Center, specifically through the System for Portable Underground Data (SPUD) Electromagnetic Transfer Function (EMTF) repository (\url{https://ds.iris.edu/spud/emtf}). The SPUD EMTF database serves as a vital global archive for MT transfer functions, providing data in the standardized Electrical Data Interchange (EDI) format. The initial dataset compiled for this research involved downloading the complete suite of available EDI files corresponding to the SAMTEX footprint, forming a massive catalog of multi-frequency electromagnetic responses sampling diverse geological domains.

\subsection{Computational Tools and Libraries}
The workflow developed for this thesis is distinctly divided into a rigorous data pre-processing and dimensionality analysis phase, followed by the computationally intensive Bayesian inversion phase.

\subsubsection{Python Ecosystem for Data Management and Geospatial Analysis}
To handle the bulk processing, filtering, and spatial analysis of the EDI files, a robust pipeline was constructed using the Python programming language. Key libraries utilized include:
\begin{itemize}
    \item \textbf{mt\_metadata:} A specialized, open-source Python package designed specifically for reading, writing, and manipulating magnetotelluric transfer function data. This library was crucial for reliably parsing the legacy EDI files, extracting the complex impedance tensor arrays, calculating derived parameters such as apparent resistivity and phase, and handling the associated metadata (e.g., station coordinates, frequencies).
    \item \textbf{Pandas and NumPy:} These foundational data science libraries were employed for structuring the extracted MT parameters into manageable DataFrames, performing vectorized mathematical operations (such as the computation of rotational invariants and Swift Skew), and executing statistical filtering to separate 1D-compatible stations from those exhibiting strong 3D geometries.
    \item \textbf{GeoPandas and Shapely:} Given the regional scope of the study, spatial context is paramount. GeoPandas was used to perform spatial joins between the MT station coordinates and digitized geological shapefiles (e.g., the McCourt 2013 tectonic map of Southern Africa). This allowed for the automated classification of each MT station into its respective tectonic unit (Archean, Limpopo, Proterozoic Belts, etc.).
    \item \textbf{Cartopy and Matplotlib:} For the visualization of the spatial distribution of stations, tectonic boundaries, and the generation of quality control (QC) pseudo-sections and sounding curves, Cartopy provided the necessary cartographic projections, while Matplotlib served as the core plotting engine.
\end{itemize}

\subsubsection{MATLAB Framework for Bayesian Inversion}
The core modeling engine of this research is a custom-developed, high-performance MATLAB codebase dedicated to 1D Trans-dimensional Hierarchical Bayesian Inversion. The architecture of this software is built upon Reversible Jump Markov Chain Monte Carlo (RJMCMC) methodologies. The code structure is highly modular, encompassing subroutines for defining prior distributions, generating synthetic forward responses (for MT, DC resistivity, or joint inversion), evaluating likelihood functions, and dynamically altering the model dimension via birth and death propositions of subsurface layers. A key feature of this custom environment is the implementation of Parallel Tempering, which manages the simultaneous execution of multiple MCMC chains across different "temperatures" to enhance global optimization and prevent the algorithm from stalling in localized probability maxima. The output of this MATLAB framework comprises massive ensembles of geoelectric models, which are subsequently analyzed to extract statistical parameters such as the median resistivity profile and its associated credible intervals (e.g., P10, P90).


\section{Methodological Advances}

The traditional approach to interpreting 1D MT data involves regularized, deterministic inversion (such as Occam's inversion), which seeks the smoothest possible model that fits the data to a predefined tolerance level. While computationally efficient, this approach imposes a subjective, user-defined structure (smoothness) on the subsurface and entirely neglects the quantification of uncertainty. To overcome these critical deficiencies, this thesis implements a suite of advanced statistical and numerical methodologies grounded in Bayesian inference.

\subsection{Trans-dimensional Inversion via RJMCMC}
In standard parametric Bayesian inversion (often driven by Metropolis-Hastings algorithms), the number of parameters—in this case, the number of horizontal layers in the 1D Earth model—must be fixed prior to the inversion. If the chosen number of layers is too low, the model will fail to capture true geological complexity (under-fitting); if it is too high, the model will fit the noise in the data, producing spurious, artificial structures (over-fitting). 

Trans-dimensional inversion elegantly solves this problem by treating the number of layers, $k$, as an unknown parameter alongside the layer resistivities and thicknesses. This is achieved through the Reversible Jump Markov Chain Monte Carlo (RJMCMC) algorithm. RJMCMC extends the standard Metropolis-Hastings acceptance criteria to accommodate jumps between parameter spaces of different dimensions. During the MCMC random walk, the algorithm proposes not only changes to the resistivity or depth of existing layers (perturbation moves) but also the creation of new layers (birth moves) or the deletion of existing ones (death moves). By integrating over a wide range of dimensions, the resulting posterior distribution naturally penalizes unnecessary complexity—a quantitative embodiment of Occam's razor. The final ensemble of models provides a probability map of resistivity at depth, where sharp interfaces are preserved if required by the data, and smooth transitions are represented by averaging across many models with varying layer boundaries.

\subsection{Hierarchical Bayesian Framework for Noise Estimation}
A persistent challenge in the inversion of real MT data is the accurate estimation of data uncertainties. Standard deviations provided in EDI files are often derived from the statistical variance of the time-series processing and may not accurately reflect the true, often non-stationary, noise profile. Errors can arise from static shift, cultural noise, or multi-dimensional geological effects that cannot be fitted by a 1D model. Incorrectly specifying the noise level in the likelihood function can severely bias the inversion, leading to either overly complex models (if noise is underestimated) or featureless, over-smoothed models (if noise is overestimated).

To address this, this research implements a Hierarchical Bayesian framework. In this approach, the standard deviation of the data errors is not fixed but is instead treated as an unknown hyper-parameter to be inferred jointly with the Earth model parameters. The noise parameter is assigned its own prior distribution (the hyper-prior) and is updated during the MCMC sampling. This allows the algorithm to dynamically scale the data uncertainties, autonomously determining the appropriate level of data fit based on the internal consistency of the observations. This hierarchical approach significantly enhances the robustness of the inversion against outliers and subjective bias in data weighting.

\subsection{Parallel Tempering}
The posterior probability space in highly non-linear, multi-modal problems like electromagnetic inversion is often exceptionally complex, characterized by numerous local maxima separated by deep valleys of low probability. A single MCMC chain can easily become trapped in one of these local maxima, failing to explore the global parameter space and resulting in an artificially narrow uncertainty estimate.

To ensure comprehensive exploration and convergence to the true global posterior distribution, this study employs Parallel Tempering (also known as replica exchange MCMC). This technique involves running multiple, interacting MCMC chains simultaneously. Each chain is assigned a "temperature" parameter, $T$. The base chain operates at $T=1$ and samples the true posterior distribution. The other chains operate at elevated temperatures ($T > 1$). Raising the temperature artificially flattens the likelihood landscape, allowing the higher-temperature chains to easily traverse probability valleys and explore disparate regions of the model space. Periodically, the algorithm proposes to swap the states of two adjacent chains. This swapping mechanism allows the base chain ($T=1$) to rapidly jump out of local minima by inheriting a globally diverse model from a hotter chain. Parallel Tempering drastically accelerates convergence and ensures that the final ensemble accurately reflects the true, potentially multi-modal, uncertainty of the geoelectric structure.

\subsection{Rigorous Dimensionality Analysis and Station Selection}
Because this study focuses on high-resolution 1D trans-dimensional inversion, it is an absolute prerequisite that the analyzed data strictly conform to the assumptions of a 1-Dimensional Earth (i.e., resistivity varies only with depth). Inverting 2D or 3D data with a 1D algorithm will inevitably produce severe artefacts.

To mitigate this, a stringent pre-processing workflow was developed to assess the dimensionality of every station in the SAMTEX dataset. This was achieved by calculating the Swift Skew, a rotationally invariant parameter derived from the impedance tensor. Swift Skew quantifies the degree of asymmetry in the tensor, which is indicative of 3D structures. For a perfectly 1D or 2D Earth, the Swift Skew is theoretically zero. In practice, a threshold of 0.2 is widely accepted in the MT community; stations exhibiting a Swift Skew less than 0.2 across the majority of their frequency band are deemed suitable for 1D approximation. A custom Python pipeline systematically sanitized the EDI files, computed the frequency-averaged Swift Skew, and classified the stations. From an initial pool of hundreds of sites, only those passing this rigorous 1D threshold were selected. This meticulously curated subset was then spatially analyzed and categorized into major tectonic units (e.g., Archean cratons vs. Proterozoic belts), ensuring that the subsequent Bayesian inversions represent genuine vertical variations in lithospheric properties rather than multidimensional artefacts.


\section{Regional Challenges and Research Gaps}

The application of Magnetotellurics in Southern Africa is fraught with challenges intrinsic to the extreme geological heterogeneity and antiquity of the region. The Southern African subcontinent is an amalgamation of distinct cratonic nuclei, most notably the Kaapvaal and Zimbabwe cratons, which are stitched together and surrounded by a complex network of mobile belts such as the Limpopo Belt, the Namaqua-Natal Province, and the Damara Orogen. These tectonic domains span billions of years of Earth history, from the early Archean to the Neoproterozoic, and have undergone multiple episodes of rifting, collision, magmatism, and metasomatism.

\subsection{Dimensionality and Static Shift}
One of the primary challenges in this environment is the pervasive presence of 2D and 3D electrical structures. The boundaries between cratons and mobile belts are often marked by near-vertical, highly conductive suture zones or major fault systems. When MT currents flow across or along these sharp lateral boundaries, they generate complex 3D electromagnetic responses that violate the assumptions of 1D modeling. Furthermore, the arid, highly resistive near-surface geology common in parts of South Africa and Namibia frequently leads to severe "static shift" — a galvanic distortion of the electric field caused by shallow, small-scale resistivity anomalies. Static shift displaces the apparent resistivity curves vertically on a log-log plot without altering the phase, leading to gross miscalculations of depth to deeper geological targets if not properly addressed or absorbed by uncertainty estimation.

\subsection{The Ambiguity of the Lithosphere-Asthenosphere Boundary (LAB)}
A central goal of deep lithospheric studies is mapping the LAB, which represents the transition from the rigid lithospheric plate to the more ductile, convecting asthenosphere below. Seismically, this boundary is often defined by a decrease in shear wave velocity. Electromagnetically, it is typically marked by a significant decrease in resistivity, as the higher temperatures and potential presence of partial melt or volatiles in the asthenosphere enhance conductivity. However, accurately defining the depth and sharpness of the LAB beneath cratons using MT is notoriously difficult. The extreme thickness and high resistivity of the cratonic roots attenuate the high-frequency signals, leaving only the longest-period data to resolve the LAB. Consequently, deterministic inversions often smear this boundary, producing gradual transitions that are highly dependent on the starting model and regularization. 

\subsection{Research Gaps in Probabilistic Imaging}
While extensive MT research has been conducted in Southern Africa, the vast majority of these studies have relied on deterministic, regularized 2D and 3D inversion codes (e.g., NLCG, Occam2D, ModEM). These landmark studies have successfully delineated large-scale cratonic roots and major conductive corridors. However, a significant research gap persists regarding the rigorous, statistical quantification of uncertainty in these geoelectric models. Deterministic models provide a single snapshot of the subsurface, masking the inherent ambiguity of the non-linear inverse problem. When comparing the subtle differences in lithospheric thickness or mid-crustal conductivity between adjacent tectonic units (e.g., the Kaapvaal vs. the Limpopo belt), it is critical to know whether the observed differences are statistically significant or merely artefacts of the inversion process. 

This thesis aims to bridge this gap by applying the trans-dimensional hierarchical Bayesian framework to a curated 1D dataset. By moving away from deterministic regularization and embracing a fully probabilistic approach, this research seeks to provide statistically robust answers to fundamental questions: What is the true range of acceptable lithospheric thicknesses beneath the Archean cratons? How sharp are the transitions in conductivity at depth? And how do these properties statistically differ across the various Proterozoic and Archean tectonic units of Southern Africa?

\section\label{sec:objectives}
\section{Objectives of the Thesis}

The overarching aim of this dissertation is to characterize the vertical geoelectric structure of the Southern African lithosphere across diverse tectonic terrains using advanced probabilistic inversion techniques, thereby providing uncertainty-quantified constraints on cratonic evolution and lithospheric architecture.

To fulfill this overarching aim, the research is guided by the following specific objectives:

\begin{enumerate}
    \item \textbf{Development of an Automated Dimensionality Analysis Workflow:} To design and implement a robust, Python-based computational pipeline capable of ingesting large volumes of regional MT data (legacy EDI files), sanitizing the data structures, and performing rigorous dimensionality analysis using rotational invariants (Swift Skew). This objective ensures that only data strictly conforming to 1D assumptions are carried forward, preventing the misinterpretation of multi-dimensional anomalies.
    
    \item \textbf{Spatial and Tectonic Categorization of MT Stations:} To integrate the filtered 1D MT stations with high-resolution regional geological and tectonic models. By utilizing geospatial tools (GeoPandas) and the McCourt (2013) tectonic map, this objective aims to precisely classify a highly curated subset of stations (e.g., 24 representative sites) into distinct geological domains: Archean Cratons, the Limpopo Belt, Paleoproterozoic Shields, Paleoproterozoic Belts, Mesoproterozoic Belts, and Neoproterozoic domains.
    
    \item \textbf{Implementation of Trans-dimensional Hierarchical Bayesian Inversion:} To adapt, optimize, and deploy a custom MATLAB-based 1D Reversible Jump Markov Chain Monte Carlo (RJMCMC) inversion algorithm. This involves configuring the Parallel Tempering parameters, defining appropriate prior bounds for resistivity and depth, and utilizing a hierarchical framework to autonomously estimate data noise, thereby ensuring objective and statistically rigorous inversion runs for the selected MT stations.
    
    \item \textbf{Uncertainty-Quantified Lithospheric Imaging:} To execute the Bayesian inversion on the selected stations and thoroughly analyze the resulting posterior probability density functions (PDFs). This objective focuses on extracting median resistivity-depth profiles along with their associated 95\% credible intervals, providing a clear visual and statistical representation of model uncertainty.
    
    \item \textbf{Comparative Tectonic Analysis:} To synthesize the inversion results and statistically compare the geoelectric signatures of the different tectonic units. This includes assessing variations in the thickness of the resistive lithospheric root, evaluating the depth to the electrically defined Lithosphere-Asthenosphere Boundary (LAB), and characterizing the presence and intensity of mid-to-lower crustal conductors across the Archean to Neoproterozoic age spectrum.
\end{enumerate}

\section{Thesis Organization}

This dissertation is systematically structured to guide the reader through the theoretical foundations, the methodological implementation, and the final tectonic interpretations. The thesis is organized into the following chapters:

\textbf{Chapter 1: Introduction (Current Chapter)} sets the stage by outlining the background and motivation for studying the Southern African lithosphere using Magnetotellurics. It introduces the key data sources, highlights the methodological advancements offered by Bayesian inference, identifies existing research gaps in the region, and defines the specific objectives and structure of the thesis.

\textbf{Chapter 2: Theoretical Foundations and Methodology} provides an in-depth exploration of the physics of Magnetotellurics and the mathematics of Bayesian inference. The first half details Maxwell's equations, EM wave propagation in the Earth, the derivation of the impedance tensor, and dimensionality analysis. The second half meticulously unpacks Bayes' Theorem, Markov Chain Monte Carlo sampling, the mechanics of the Reversible Jump (Trans-dimensional) algorithm for dynamically altering model layers, and the principles of Parallel Tempering and Hierarchical noise estimation.

\textbf{Chapter 3: Data Processing and Spatial Analysis} details the practical application of Python to the SAMTEX dataset. It chronicles the transition from raw EDI files to clean, format-compliant data arrays. This chapter heavily focuses on the calculation of Swift Skew, the rigorous filtering process to isolate 1D-compatible stations, and the geospatial methodologies utilized to intersect these stations with the McCourt (2013) tectonic map, culminating in the selection of the final 24 high-quality stations utilized for Bayesian inversion.

\textbf{Chapter 4: Results of the Trans-dimensional Bayesian Inversion} presents the core findings of the computational modeling. For each of the selected 24 stations, this chapter displays the MCMC chain convergence diagnostics, the fit to the observed data, and the final posterior probability density maps of the 1D resistivity structure. The results are grouped by tectonic unit, providing a detailed, statistically robust view of the crust and upper mantle architecture beneath the Archean, Limpopo, and various Proterozoic domains.

\textbf{Chapter 5: Discussion and Tectonic Synthesis} interprets the inversion results in a broader geological context. It compares the thickness of the resistive lithosphere and the depth to the LAB across the different tectonic terranes. This chapter discusses the implications of these findings for the stabilization of the Kaapvaal Craton, the nature of the boundaries between the craton and the mobile belts, and the origin of mid-crustal conductors. The chapter concludes by summarizing the contributions of this research and proposing avenues for future trans-dimensional studies.

\newpage
\section{Basic Terminologies and Definitions}

To ensure clarity and consistency throughout this dissertation, the following subsection provides detailed definitions and contextual explanations of the fundamental physical concepts, mathematical parameters, and computational algorithms that form the core of this research.

\subsection{1.1.1 Magnetotellurics (MT)}
Magnetotellurics is a passive, frequency-domain electromagnetic sounding technique used to image the electrical resistivity structure of the Earth's subsurface. Unlike controlled-source methods that require a transmitter, MT relies on naturally occurring, time-varying electromagnetic fields generated in the Earth's ionosphere and magnetosphere (from solar wind interactions) and global lightning activity. These external fields induce secondary electric currents (telluric currents) within the Earth. By simultaneously measuring the orthogonal components of the fluctuating magnetic field ($H_x, H_y$) and the induced electric field ($E_x, E_y$) at the surface over a wide range of frequencies, geophysicists can determine the subsurface conductivity distribution from shallow depths (tens of meters) down to the deep upper mantle (hundreds of kilometers).

\subsection{1.1.2 Electrical Resistivity ($\rho$)}
Electrical resistivity, denoted by $\rho$, is an intrinsic physical property of a material that quantifies how strongly it opposes the flow of electric current. In the International System of Units (SI), it is measured in Ohm-meters ($\Omega \cdot m$). It is the reciprocal of electrical conductivity ($\sigma$, measured in Siemens per meter, $S/m$). In the Earth's lithosphere, silicate minerals are generally excellent insulators, meaning cold, dry, crystalline rock has very high resistivity ($>10,000 \ \Omega \cdot m$). However, resistivity drops exponentially with increasing temperature. Furthermore, the presence of even small interconnected fractions of conductive phases—such as saline fluids, partial melts, graphite films, or metallic sulfides—can drastically lower the bulk resistivity of the rock. Thus, resistivity is an exceptionally sensitive proxy for the thermal and compositional state of the deep Earth.

\subsection{1.1.3 Impedance Tensor ($Z$)}
The fundamental transfer function derived from MT measurements is the complex impedance tensor, denoted as $\mathbf{Z}$. It linearly relates the horizontal components of the electric field ($\mathbf{E}$) to the orthogonal horizontal components of the magnetic field ($\mathbf{H}$) at a given angular frequency ($\omega$). Mathematically, it is expressed as $\mathbf{E} = \mathbf{Z}\mathbf{H}$, or in matrix form:
\begin{equation}
\begin{bmatrix} E_x \\ E_y \end{bmatrix} = \begin{bmatrix} Z_{xx} & Z_{xy} \\ Z_{yx} & Z_{yy} \end{bmatrix} \begin{bmatrix} H_x \\ H_y \end{bmatrix}
\end{equation}
The elements of the tensor ($Z_{xx}, Z_{xy}, Z_{yx}, Z_{yy}$) are complex numbers containing magnitude and phase information. The structure and symmetry of this tensor depend entirely on the dimensionality of the underlying subsurface. For a purely 1D Earth (horizontally layered), the diagonal elements ($Z_{xx}, Z_{yy}$) are zero, and the off-diagonal elements are equal in magnitude but opposite in sign ($Z_{xy} = -Z_{yx}$).

\subsection{1.1.4 Apparent Resistivity ($\rho_a$) and Phase ($\phi$)}
Because the Earth is rarely a uniform half-space, the resistivity calculated directly from the impedance tensor is an "apparent" value, representing a volume-averaged resistivity of the heterogeneous subsurface sampled by the EM wave at a specific frequency. For a specific component (e.g., $xy$), the apparent resistivity is calculated as:
\begin{equation}
\rho_{a,xy}(\omega) = \frac{1}{\mu_0 \omega} |Z_{xy}(\omega)|^2
\end{equation}
where $\mu_0$ is the magnetic permeability of free space. The Phase ($\phi$) represents the time lag (or phase shift) between the measured electric and magnetic fields. It is calculated from the ratio of the imaginary to the real part of the complex impedance:
\begin{equation}
\phi_{xy} = \tan^{-1} \left( \frac{\text{Im}(Z_{xy})}{\text{Re}(Z_{xy})} \right)
\end{equation}
The phase is particularly sensitive to vertical gradients in true resistivity; a phase greater than 45 degrees indicates a transition into a more conductive medium at depth, while a phase less than 45 degrees indicates a transition into a more resistive medium.

\subsection{1.1.5 Skin Depth ($\delta$)}
Electromagnetic fields diffuse into the Earth and attenuate exponentially with depth. The "Skin Depth" is defined as the depth at which the amplitude of an incident plane EM wave decays to $1/e$ (approximately 37\%) of its value at the surface. It serves as a rough proxy for the depth of investigation for a specific frequency. The skin depth ($\delta$, in meters) is given by:
\begin{equation}
\delta \approx 503 \sqrt{\frac{\rho}{f}}
\end{equation}
where $\rho$ is the average resistivity of the host rock ($\Omega \cdot m$) and $f$ is the frequency of the signal in Hertz. This fundamental equation illustrates that low-frequency (long-period) signals penetrate much deeper into the Earth than high-frequency signals, and that EM waves penetrate deeper into highly resistive environments (like cratons) than into conductive environments.

\subsection{1.1.6 Swift Skew}
Swift Skew is a rotational invariant parameter calculated from the elements of the impedance tensor, used to quantify the degree to which the local geoelectric structure deviates from 1D or 2D geometry into 3D complexity. It essentially measures the asymmetry of the tensor. It is defined as:
\begin{equation}
\text{Skew} = \frac{|Z_{xx} + Z_{yy}|}{|Z_{xy} - Z_{yx}|}
\end{equation}
In an ideal 1D or perfectly 2D Earth, the Swift Skew is exactly zero. As the structural complexity increases (e.g., near corners of geological bodies or intersecting fault zones), the skew value rises. In MT processing, a heuristic threshold is often applied: if the Swift Skew is less than 0.2 across a broad frequency range, the site is generally considered suitable for 1D or 2D approximation. Values significantly higher than 0.2 necessitate fully 3D inversion methodologies.

\subsection{1.1.7 Forward Modeling}
Forward modeling is the mathematical process of predicting the theoretical data (e.g., apparent resistivity and phase curves) that would be observed at the surface given a known or hypothesized model of the subsurface. In the context of a 1D Earth, forward modeling involves solving Maxwell's equations analytically for a sequence of horizontal layers, each defined by a specific thickness and uniform resistivity. An iterative inversion algorithm relies heavily on repeated forward modeling: it proposes a model, calculates the forward response, compares this synthetic response to the actual measured field data (calculating the misfit), and then updates the model parameters to reduce this misfit.

\subsection{1.1.8 Bayesian Inference}
Bayesian inference is a method of statistical reasoning grounded in Bayes' Theorem, which updates the probability of a hypothesis as more evidence or information becomes available. In geophysical inversion, the "hypothesis" is the set of Earth model parameters ($\mathbf{m}$), and the "evidence" is the observed MT data ($\mathbf{d}$). Bayes' Theorem is expressed as:
\begin{equation}
P(\mathbf{m}|\mathbf{d}) = \frac{P(\mathbf{d}|\mathbf{m}) P(\mathbf{m})}{P(\mathbf{d})}
\end{equation}
Here, $P(\mathbf{m}|\mathbf{d})$ is the Posterior probability (the goal of the inversion—the probability of the model given the data). $P(\mathbf{d}|\mathbf{m})$ is the Likelihood function (how well the forward model fits the data). $P(\mathbf{m})$ is the Prior probability (what is known about the geology before seeing the data, such as reasonable upper and lower bounds for resistivity). $P(\mathbf{d})$ is the Evidence, a normalizing constant. Unlike deterministic methods that seek a single maximum of the Posterior, Bayesian methods seek to map the entire Posterior distribution.

\subsection{1.1.9 Reversible Jump Markov Chain Monte Carlo (RJMCMC)}
Markov Chain Monte Carlo (MCMC) is a class of algorithms for sampling from complex probability distributions (like the Bayesian Posterior) by constructing a Markov chain that has the desired distribution as its equilibrium distribution. Reversible Jump MCMC (RJMCMC) is an advanced extension developed by Peter Green in 1995 that allows the Markov chain to "jump" between models of different dimensionalities. In 1D MT inversion, the dimension is the number of geological layers. An RJMCMC algorithm will propose standard MCMC moves (perturbing the resistivity or depth of an existing layer) alongside trans-dimensional moves: proposing to split an existing layer into two (a "birth" move, increasing dimension) or combining two adjacent layers into one (a "death" move, decreasing dimension). This allows the algorithm to automatically discover the optimal level of complexity supported by the data.

\subsection{1.1.10 Parallel Tempering}
Parallel Tempering (also known as Replica Exchange MCMC) is an advanced sampling technique designed to prevent MCMC chains from becoming permanently trapped in local probability maxima, a common issue in highly non-linear problems. Instead of running a single MCMC chain, Parallel Tempering runs an ensemble of $N$ chains concurrently. Each chain is assigned a "temperature" parameter $T_i$. The first chain runs at $T_1 = 1$ and targets the true posterior distribution. The other chains run at elevated temperatures ($T_2 < T_3 < \dots < T_N$). Raising the temperature artificially flattens the likelihood function, allowing the "hot" chains to freely explore the entire parameter space and escape deep local minima. Periodically, the algorithm proposes to swap the current states of two adjacent chains based on a specific acceptance probability. This allows the primary chain ($T=1$) to inherit widely differing, yet globally plausible models discovered by the hot chains, drastically improving the mixing and global convergence of the algorithm.

